\documentclass[aps,pre,reprint,superscriptaddress,nofootinbib]{revtex4-2}

\usepackage{amsmath}
\usepackage{amssymb}
\usepackage{bm}
\usepackage{graphicx}
\usepackage[colorlinks=true,allcolors=blue]{hyperref}

\newcommand{\DKL}{D_{\mathrm{KL}}}
\newcommand{\Acal}{\mathcal{A}}

\begin{document}

\title{Controlled numerical four-level counterexample to triplet screening of the Markovian Mpemba effect}

\author{Xiaojian Zeng}
\email{xiaojianzeng@foxmail.com}
\affiliation{PBC School of Finance, Tsinghua University, Beijing 100190, China}

\date{September 12, 2026}

\begin{abstract}
Local rate tests derived for three-level subsystems have been proposed as a scalable
screen for the Markovian Mpemba effect in larger detailed-balance systems.  We test
the finite-temperature proposal in the minimal four-level setting and give a simple
controlled-numerical witness.  In units where the bath inverse temperature and a
common rate scale are one, the energies are
$(0,1/200,7/20,14)$ and the six downward rates are
$(1/10,1,13,99/100,1,3/4)$.  Every one of the four triplets strictly fails its
applicable published necessary-condition predicate, while the slow nonstationary
mode is simple.  Nevertheless, Gibbs preparations at inverse temperatures
$1/50$ and $1/20$ have normalized slow-mode amplitudes
$0.0628307488472$ and $0.0630538586374$, respectively.  Thus the initially hotter
state has the smaller late-time amplitude, and direct relative-entropy evolution
reverses the initial distance ordering.  Separate generator constructions,
100- and 180-digit calculations, residual tests, and direct dynamics reproduce the
effect.  The result is a reproducible computational counterexample at controlled
numerical scope, not a directed-rounding proof or a classification of the heating
branch or of general $N$.
\end{abstract}

\maketitle

\section{Introduction}

The Mpemba effect denotes anomalous relaxation in which a preparation initially
farther from a common equilibrium state becomes closer at a later time.  Modern
formulations separate this phenomenon from the material-specific question of
freezing water and place it in the spectral theory of nonequilibrium relaxation
\cite{Teza2026}.  For finite Markov systems obeying detailed balance, the overlap of
a thermal preparation with the slowest nonstationary mode gives a particularly
clean late-time criterion \cite{LuRaz2017,Klich2019}.

This criterion is simple once the full generator has been diagonalized, but it does
not by itself provide an inexpensive way to screen a large parameter space.  Avitan,
Factor, and Gelbwaser-Klimovsky recently derived elementary necessary conditions for
a three-level system and proposed applying them to every triplet of an $N$-level
generator \cite{Avitan2026}.  They proved an exclusion statement at zero bath
temperature under stronger inequalities and reported finite-temperature random
tests, while the implication from triplet exclusion to global exclusion at finite
temperature remained conjectural.  A subsequent study found four-level Mpemba
behavior even when no isolated triplet itself exhibits the effect
\cite{GelbwaserKlimovsky2026}; that premise is weaker than failure of every published
necessary-condition predicate.  Computational screening has also been explored in
other Markovian settings \cite{Amorim2023}.

Here we address the stronger, unresolved premise directly.  We exhibit one
four-level, finite-positive-temperature, detailed-balance generator for which all
four applicable triplet predicates fail strictly, yet two cooling preparations have
the reversed slow-mode ordering.  The exact input parameters are rational, but the
Boltzmann factors and eigensystem are evaluated numerically.  We therefore state the
result at controlled-numerical scope and retain explicit margins and reproducibility
checks rather than presenting it as an exact disproof.

\section{Model and late-time criterion}

Consider four strictly ordered levels $E_1<E_2<E_3<E_4$ and a population column
vector $\bm p(t)$.  Its autonomous dynamics is
\begin{equation}
 \begin{aligned}
  \dot{\bm p}&=M\bm p,
  & M_{ij}&=a_{ij}>0\quad(i\ne j),\\
  M_{jj}&=-\sum_{i\ne j}a_{ij}.&&
 \end{aligned}
  \label{eq:master}
\end{equation}
where $a_{ij}$ is the rate for $j\to i$.  A fixed bath at inverse temperature
$\beta_b$ imposes detailed balance,
\begin{equation}
 a_{ij}\pi_j=a_{ji}\pi_i,
 \qquad
 \pi_i=\frac{e^{-\beta_b E_i}}{Z(\beta_b)}.
 \label{eq:db}
\end{equation}
For $i<j$, we take the six downward rates $d_{ij}=a_{ij}$ as independent and set
$a_{ji}=d_{ij}e^{-\beta_b(E_j-E_i)}$.  Strict positivity makes the finite chain
irreducible.  A Gibbs preparation at inverse temperature $\beta$ is denoted by
$\bm g(\beta)$, with $g_i(\beta)=e^{-\beta E_i}/Z(\beta)$; changing $\beta$ changes
only the initial state, not the bath or the generator.

Let $D=\operatorname{diag}(\bm\pi)$.  Detailed balance makes
$S=D^{-1/2}MD^{1/2}$ real symmetric.  Order the eigenvalues as
$0=\lambda_1>\lambda_2>\lambda_3\geq\lambda_4$ and assume that $\lambda_2$ is
simple.  If $\bm u_2$ is a real unit eigenvector of $S$, define
\begin{equation}
 \begin{aligned}
 \bm v_2&=D^{1/2}\bm u_2,
 &\bm\ell_2&=D^{-1/2}\bm u_2,\\
 b_2(\beta)&=\bm\ell_2^{\mathsf T}[\bm g(\beta)-\bm\pi].&&
 \end{aligned}
 \label{eq:mode}
\end{equation}
and the normalization-independent preparation amplitude
$\Acal(\beta)=|b_2(\beta)|$.  With this fixed normalization,
\begin{equation}
 \DKL[\bm p(t;\beta)\Vert\bm\pi]
 =\frac{b_2(\beta)^2}{2}e^{2\lambda_2 t}
  +o(e^{2\lambda_2 t}).
 \label{eq:klasymp}
\end{equation}
Consequently, for two cooling preparations
$0<\beta_h<\beta_c<\beta_b$, the strict reversal
$\Acal(\beta_h)<\Acal(\beta_c)$ guarantees that the hotter state is closer in
relative entropy at sufficiently late times, even if it starts farther away.

\section{Triplet screening predicate}

For an ordered triplet $i<j<k$, retain the off-diagonal rates of the full
four-level generator and define
\begin{align}
 q&=e^{-\beta_b(E_j-E_i)}, &
 r&=e^{-\beta_b(E_k-E_j)},\nonumber\\
 \kappa&=\frac{a_{ik}+a_{jk}}{a_{ij}(1+q)},\label{eq:kappa}\\
 F_{>}&=a_{ik}\left(1+\frac{qr}{2}\right)
       -a_{jk}\left(1+\frac{r}{2}\right)\nonumber\\
      &\hspace{2.4cm}-a_{ij}(1-q),\label{eq:fhigh}\\
 F_{<}&=a_{ik}q-a_{jk}.\label{eq:flow}
\end{align}
The published three-level necessary condition is $F_{>}>0$ when $\kappa\geq1$
and $F_{<}>0$ when $\kappa\leq1$ \cite{Avitan2026}.  At $\kappa=1$ the two signs
agree; none of the triplets below lies on that boundary.  The proposed global
screen excludes a larger system when every triplet fails its applicable strict
condition.  Importantly, this test uses the original internal edge rates; it is not
a principal-submatrix generator with external escape terms, an effective-rate
model, or an adiabatic elimination.

\section{Four-level witness}

Set $\beta_b=1$ by measuring energies in units of $\beta_b^{-1}$ and measure rates
in a common inverse-time unit $\gamma$.  The witness is
\begin{align}
 (E_1,E_2,E_3,E_4)&=\left(0,\frac{1}{200},\frac{7}{20},14\right),\label{eq:energies}\\
 (d_{12},d_{13},d_{23},d_{14},d_{24},d_{34})
 &=\gamma\left(\frac{1}{10},1,13,\frac{99}{100},1,\frac{3}{4}\right).
 \label{eq:rates}
\end{align}
Table~\ref{tab:triplets} evaluates Eqs.~\eqref{eq:kappa}--\eqref{eq:flow}.
For compact comparison we also give
$m_{ijk}=-F_{\mathrm{app}}/(a_{ik}+a_{jk}+a_{ij})$.  Every applicable value is
negative, and the smallest normalized exclusion margin is
$2.8325184\times10^{-3}$.

\begin{table}[t]
\caption{Triplet tests for the witness.  $F_{\mathrm{app}}$ is $F_{>}$ or $F_{<}$
according to the listed branch.  Positive $m_{ijk}$ means that the applicable
strict necessary condition is excluded.}
\label{tab:triplets}
\begin{ruledtabular}
\small
\begin{tabular}{ccccc}
Triplet & $\kappa$ & branch & $F_{\mathrm{app}}/\gamma$ & $m_{ijk}$\\
\hline
$123$ & $70.1750$ & $>$ & $-16.25159$ & $1.152595$\\
$124$ & $9.97487$ & $>$ & $-0.0104988$ & $0.00502333$\\
$134$ & $1.02071$ & $>$ & $-0.0553119$ & $0.0201868$\\
$234$ & $0.0788045$ & $<$ & $-0.0417796$ & $0.00283252$\\
\end{tabular}
\end{ruledtabular}
\end{table}

The relaxation spectrum, ordered from stationary to fastest, is
\begin{equation}
 \begin{split}
 \frac{\bm\lambda}{\gamma}={}&
 (0,-1.25273315932984,\\
 &-2.74000084438500,-22.8583224729113).
 \end{split}
 \label{eq:spectrum}
\end{equation}
Thus the slow mode is simple.  Its separation from the next mode, normalized by
the full spectral width, is
$(\lambda_2-\lambda_3)/(-\lambda_4)=0.0650646033548$.

Now take two Gibbs preparations on the cooling branch,
$\beta_h/\beta_b=1/50$ and $\beta_c/\beta_b=1/20$.  Equation~\eqref{eq:mode}
gives
\begin{align}
 \Acal(\beta_h)&=0.0628307488471692912,\nonumber\\
 \Acal(\beta_c)&=0.0630538586374105295,\nonumber\\
 \Acal(\beta_c)-\Acal(\beta_h)
 &=2.2310979024123833\times10^{-4}>0.
 \label{eq:amplitudes}
\end{align}
This is the ordering forbidden by the proposed cooling screen.  The hotter state
does start farther from equilibrium:
$\DKL[\bm g(\beta_h)\Vert\bm\pi]=2.53001045904$, whereas
$\DKL[\bm g(\beta_c)\Vert\bm\pi]=1.74010206475$.  Direct propagation with the
same fixed generator gives $D_h-D_c=1.6527113\times10^{-10}$ at
$\gamma t=6$ and $-2.4701607\times10^{-14}$ at $\gamma t=8$, so the checked
distance ordering changes between those times.  The small late-time difference is
not the sole discriminator: the resolved amplitude inequality in
Eq.~\eqref{eq:amplitudes} independently fixes the eventual ordering through
Eq.~\eqref{eq:klasymp}.

\begin{figure*}[t]
  \includegraphics[width=0.98\textwidth]{manuscript/fig-counterexample.pdf}
  \caption{Three views of the same frozen numerical witness.  (a) All four
  normalized exclusion margins are positive.  (b) The normalized slow-mode
  amplitude increases from the hotter preparation ($\beta/\beta_b=0.02$) to the
  cooler one ($0.05$), reversing the proposed cooling monotonicity.  (c) Relative
  entropy initially favors the cooler state, while the directly calculated
  difference $D_h-D_c$ becomes negative at the checked late times.  Markers are
  archived high-precision outputs and connecting lines are guides to the eye.  The
  plotting code was generated with substantive assistance from OpenAI Codex; it
  deterministically reads the archived machine-readable verifier output without
  manual data entry.}
  \label{fig:witness}
\end{figure*}

\section{Discovery and numerical control}

The mechanism was found by a deterministic finite scan followed by constrained
local optimization.  That finite search is discovery provenance rather than a global
certificate.  After identifying a boundary mechanism---a nearly degenerate lowest
pair together with a widely separated top level---we rationalized the parameters and
froze Eqs.~\eqref{eq:energies} and \eqref{eq:rates} before the high-precision
calculation.

For verification, the generator was assembled in two separately coded conventions:
an ascending-energy column generator following Eq.~\eqref{eq:master}, and a
descending-energy row generator mapped back by reversal and transposition.  The
mapped matrices agree.  The detailed-balance similarity was diagonalized at 100
and 180 decimal digits, and the amplitudes, spectral gap, triplet predicates, and
matrix-exponential dynamics were recalculated at both precisions.  The decisive
quantities agree beyond 90 decimal places.  At 180 digits the largest reported
matrix or modal residual is below $2\times10^{-180}$, compared with the
$2.23\times10^{-4}$ amplitude difference and the $2.83\times10^{-3}$ smallest
normalized predicate margin.  Probability conservation, stationarity, detailed
balance, positivity along the checked trajectories, left/right eigenpair equations,
and biorthogonal normalization were tested explicitly.

These checks control ordinary floating-point and representation errors very strongly,
but arbitrary precision is not directed rounding.  We have not enclosed the Gibbs
factors or the isolated eigenpair with intervals and do not attach a theorem-level
sign certificate to the result.  The exact-rational input and large stable margins
make the witness reproducible as a computational result; stronger exact or
computer-assisted language would require a new enclosure argument.

\subsection{AI-use disclosure for the computational workflow}

The user designated the source conversation and required a staged verification
workflow.  The topic originated in an AI-generated batch on that designated page,
which the user also supplied as an attachment; a local AI supervisor then selected
this candidate and formalized its scope before research.  OpenAI Codex using the
GPT-6 Astra and GPT-5.6 Sol model families (accessed September 2026) provided
substantive assistance with
literature synthesis, numerical-search design, code generation and debugging,
verification planning, figure-code generation, and scientific interpretation.  The
AI-assisted output was constrained by a frozen model statement and checked through
the two generator representations, precision scaling, residual tests, and direct
dynamics described above.  No AI system is an author.  Before submission, the named
author must independently inspect the cited sources, code, outputs, figures, and
claims, revise them as needed, and accept full responsibility; that human review is
pending for this draft.

\section{Physical interpretation and limitations}

The witness is not tied to an unphysical sign choice.  In the chosen rate unit, take
a common Arrhenius attempt rate $\Gamma=14\gamma$.  For every edge $i<j$, choose a
symmetric barrier satisfying
\begin{equation}
 \beta_b B_{ij}=\beta_b E_j+\ln(14\gamma/d_{ij}).
 \label{eq:barrier}
\end{equation}
Because every $d_{ij}<14\gamma$, the barrier lies above the upper level.  The
downward rate is $\Gamma e^{-\beta_b(B_{ij}-E_j)}=d_{ij}$ and the reverse rate is
$d_{ij}e^{-\beta_b(E_j-E_i)}$, exactly as in Eq.~\eqref{eq:db}.  This supplies one
standard ideal-bath realization of the population generator; it is not a claim of a
unique or universal microscopic spectral density.

The strict predicate margins, nonzero spectral separation, and strict amplitude
difference also imply continuity under sufficiently small parameter perturbations,
so the effect is not confined to a tuned equality.  We do not quantify that
neighborhood here.  The large scaled top energy, $\beta_bE_4=14$, lies outside the
compact energy boxes used in the reported finite-temperature random tests
\cite{Avitan2026,GelbwaserKlimovsky2026}; this helps explain why a finite no-hit scan
does not establish the proposed implication.

The scope is deliberately narrow.  We show one cooling pair for one four-state
classical population generator with a simple slow mode.  We do not classify inverse
heating, degenerate slow subspaces, systems with $N>4$, coherence, memory, a changing
bath, transient crossings of unrelated observables, or microscopic universality.
Nor does a bounded literature search establish absolute priority.  Within these
limits, the calculation demonstrates the physical point: local three-level
necessary-condition screening can miss collective four-level anomalous cooling.

\section{Conclusion}

A finite-temperature four-level detailed-balance generator can fail every applicable
three-level necessary-condition predicate and still display a Markovian Mpemba
cooling reversal.  The explicit rational fixture, resolved slow mode, stable amplitude
margin, and direct entropy overtaking make this conclusion reproducible at controlled
numerical scope.  The example therefore invalidates the proposed universal cooling
screen at that scope, while leaving the heating branch, general $N$, and exact
computer-assisted certification open.

\section*{Data availability}

The machine-readable high-precision output, the independent-convention verifier,
the plotting script, and environment requirements accompany this manuscript.  The
verifier reconstructs the generator from the exact
parameters in Eqs.~\eqref{eq:energies} and \eqref{eq:rates}; no downloaded paper or
private runtime log is required to reproduce the numerical tables and figure.

\begin{acknowledgments}
OpenAI Codex also assisted substantively with organization and drafting of this
manuscript.  The scope and pending author-review requirement are stated in the
AI-use disclosure above.
\end{acknowledgments}

\bibliography{manuscript/references}

\end{document}
